\section{课程定位与主问题}

本讲由 Weizhu Chen（Microsoft）主讲，主题不是 ``如何再造一个 Agent 框架''，而是 ``在真实工业环境里，Agentic 模型到底如何被训练出来''。演讲一开始就明确将重点放在 \textbf{post-training}，尤其是 reinforcement learning（RL）驱动的 agent 能力构建，而不是推理框架封装或 prompt 工程技巧。

从字幕可以看到，讲者反复强调一个实践判断：\textit{``if you want to improve model quality, you need end-to-end training.''} 这里的 end-to-end 不是指单一网络结构，而是指从数据定义、grader 设计、rollout 生成、policy 更新到产品反馈回流的闭环链路。这一视角决定了整节课的技术重心。

\begin{importantbox}{本讲最核心的一句话}
Agentic 模型训练的瓶颈不在 ``会不会调用工具''，而在 ``能否稳定完成目标导向的多轮决策并被可执行地评估''。因此，真正的工程工作量主要在数据、grader 与系统效率三处。
\end{importantbox}

讲者以 coding agent 作为主例，不是因为 coding 是唯一场景，而是因为该场景同时具备以下条件：任务目标清晰、环境反馈可程序化、失败可重复、评估可自动化。对课程学习者而言，这提供了一个高信噪比的观测窗口。

\begin{knowledgebox}{从研究视角到产品视角的转换}
在研究论文里，我们常把改进写成 ``算法提升''；在产品中，改进往往来自 ``评价体系重构 + 数据再配比 + 系统吞吐优化''。这也是讲者为何将大量时间投入在数据和 grader 上，而非只讲 RL loss。
\end{knowledgebox}

\begin{table}[H]
\centering
\begin{tabular}{p{2.6cm}p{4.2cm}p{4.7cm}}
\toprule
\textbf{问题层} & \textbf{学术常见表述} & \textbf{本讲对应的工业表述} \\
\midrule
目标定义 & maximize reward & 明确定义产品可接受行为与交互完成标准 \\
数据来源 & supervised / RL data & 可验证样本、合成样本、在线交互回流样本 \\
评估指标 & pass@1, reward score & 质量、时延、格式、稳定性、多轮一致性 \\
系统成本 & training FLOPs & sampler 吞吐、GPU 利用率、失败恢复时间 \\
\bottomrule
\end{tabular}
\caption{本讲将 Agent 训练问题从论文范式切换到生产范式}
\end{table}

\subsection{本章小结}
本讲的定位是 ``工业级 Agentic 模型训练实践''。如果只把它当作 RL 算法课，会漏掉真正决定质量上限的三个变量：数据构建、grader 设计和大规模系统效率。

